% Copyright (c) 2008-2009 solvethis
% Copyright (c) 2010-2016,2018-2019,2021 Casper Ti. Vector
% Copyright (c) 2021 Kurapica
% Copyright (c) 2021 iofu728
% Copyright (c) 2025 hylz
%
% 当前overleaf 版本字体和格式均符合2021硕士学位论文要求
%
% 该版本遵循北京大学研究生学位论文写作指南2019年v2版本、北京大学硕士研究生学位论文word模板，
% 在CasperVector/pkuthss模板的基础上针对硕士研究生学位论文格式的要求进行相应修改，
% 遵循 LaTeX Project Public License 和 知识共享 署名 - 非商业性 - 相同方式共享 4.0 国际协议
% 如有任何疑问请在github/iofu728/pkuthss上提问或联系作者iofu728。
%
% hylz使用Gemini 3 Pro (Preview)基于github/iofu728/pkuthss v1.8.3进行了封面格式的修改，使其适用于上海创智学院的博士研究生开题报告。
%
% 此处请保留 ugly 参数，参考文献管理采用GB/T 7714-2015标准。
% 此处使用顺序编码制，如使用著者-出版年制则更改为b7714-2015av。
% bib引用使用基本与日常学术论文写作相同，部分略有不同，
% 参考https://github.com/hushidong/biblatex-gb7714-2015/blob/master/example/cls-beamer.tex
% 关键：强制 ctex 使用本仓库的字体集定义（ctex-fontset-siithss.def），
% 避免 TeX Live 自带 ctex-fontset-mac.def 去寻找 STSong.otf/STKaiti.otf 等不存在的字体文件。
\documentclass[UTF8,fontset=siithss,nocolorlinks,ugly]{siithss}

% \usepackage[backend=biber,bibstyle=gb7714-2015,citestyle=gb7714-2015]{biblatex}
% \usepackage[backend=biber,style=cvpr]{biblatex}
\usepackage{float}
\usepackage{booktabs}
\usepackage{makecell}
\usepackage[nottoc,notlot,notlof]{tocbibind} % 参考文献加入目录
\usepackage{bm}
\usepackage{wrapfig}
\usepackage{colortbl}
%\usepackage[dvipsnames,table]{xcolor} % colors
\usepackage{multirow}
\usepackage{amsfonts} % blackboard math symbols
\usepackage{nicefrac} % compact symbols for 1/2, etc.
\usepackage{microtype} % microtypography

% 图表等引用汉字与数字之间不加空格
\newcommand{\Tref}[1]{表\ref{#1}}
\newcommand{\tref}[1]{表\ref{#1}}
\newcommand{\Eref}[1]{式(\ref{#1})}
\newcommand{\Fref}[1]{图\ref{#1}}
\newcommand{\Sref}[1]{第\ref{#1}节}
\newcommand{\Aref}[1]{算法\ref{#1}}
\newcommand{\eref}[1]{式(\ref{#1})}
\newcommand{\fref}[1]{图\ref{#1}}
\newcommand{\sref}[1]{第\ref{#1}节}

\newcommand{\parhead}[1]{\vskip 2pt \noindent{\bfseries #1\  }}

\newcommand{\ev}{{e}}
\newcommand{\fr}{{I}}
\newcommand{\mev}{{E}}
\newcommand{\mfr}{{L}}
\newcommand{\mdi}{{D}}
\newcommand{\mcf}{{C}}
\newcommand{\vect}[1]{{\bm{#1}}} %%Vector
\newcommand{\vp}{\vect{p}}
\newcommand{\ppix}{{\vp}}
\newcommand{\tcon}{{t}}
\newcommand{\tdis}{{\tau}}
\newcommand{\texp}{{T}}
\newcommand{\thres}{{c}}
\newcommand{\opt}[1]{\mathcal{#1}} %%Operator
\newcommand{\oF}{\opt{F}}
\newcommand{\oL}{\opt{L}}

\usepackage{pifont}
\newcommand{\cmark}{\text{\ding{51}}}
\newcommand{\xmark}{\text{\ding{55}}}

\definecolor{tabfirst}{rgb}{1, 0.7, 0.7} % red
\definecolor{tabsecond}{rgb}{1, 0.85, 0.7} % orange
\definecolor{tabthird}{rgb}{1, 1, 0.7} % yellow
\definecolor{lgg}{rgb}{0.75, 0.95, 0.65} % 定义一个名为 lightgrassgreen 的颜色

%\setlength{\bibitemsep}{3bp}
%\renewcommand*{\bibfont}{\zihao{5}\linespread{1.27}\selectfont}

\siithssinfo{
	ctitle = {基于\LaTeX 的开题报告撰写模板},
	etitle = {A \LaTeX ~Template for Doctoral Research Proposal},
	cauthor = {创小智}, eauthor = {Xiaozhi Chuang},
	studentid = {20251218},
	date = {\zhdigits{2025} 年 \zhnumber{12} 月},
	% 根据报告实际日期修改
	cmajor = {计算机应用技术}, emajor = {Computer Application Technology},
	type = {thesis}, % thesis 或 practice, 对应“学位论文”或“实践成果”
	mentorleader = {龙创创},
	mentormembers = {龙智智~~龙慧慧},
	ckeywords = {关键词1，关键词2，关键词3},
	ekeywords = {Key Words 1, Key Words 2, Key Words 3}
}
%\addbibresource{ref.bib}


\begin{document}
	\frontmatter
	\pagestyle{empty}
	\maketitle

	\chapter*{摘要}

撰写高质量的开题报告是博士研究生培养过程中的重要环节。一份规范、美观的开题报告不仅能够清晰地展示研究思路和计划，还能体现研究者的学术素养。随着计算机排版技术的发展，\LaTeX 因其强大的公式排版能力和优秀的文档管理功能，已成为学术界广泛使用的文档撰写工具。然而，上海创智学院的博士生开题报告目前只有基于Word的模板，缺乏专门针对\LaTeX 的开题报告模板，导致许多研究生在撰写过程中面临诸多挑战。

本课题旨在设计并实现一套适用于上海创智学院博士生的开题报告\LaTeX 模板。通过分析现有开题报告的结构和格式要求，结合\LaTeX 的排版优势，开发出一套功能完善、易于使用的模板。围绕这一研究方向，本课题从三个方面展开，分别探索了工作一、工作二和工作三的研究方法与技术路线，取得了初步的实验成果。在未来工作中，本课题拟从方面一、方面二、方面三继续深化研究，进一步完善模板功能，提高用户体验。

\vspace{3em}
\noindent{\bfseries 关键词：}\LaTeX，开题报告，模板设计，学术写作，排版自动化

\clearpage % 摘要

    \cleardoublepage
	\pagestyle{plain}
	\setcounter{page}{0}
	\pagenumbering{Roman}
	\tableofcontents

	\mainmatter
	\markboth{选题背景和意义}{选题背景和意义}
\chapter{选题背景和意义}%\thispagestyle{checklist}
\label{chap:background}

本章主要介绍本课题的研究背景、研究意义、国内外研究现状以及本文的研究目标和内容安排。

\section{博士生开题报告的背景与意义}

博士生开题报告是博士学位论文工作的重要环节，标志着博士生正式进入学位论文研究阶段。通过开题报告，博士生需要系统梳理相关领域的研究现状，明确研究问题，论证选题的科学价值和现实意义。这一过程不仅有助于理清研究思路，规避潜在的研究风险，还能通过专家组的评审与建议，进一步完善研究方案，确保后续研究工作的顺利开展与高质量完成。

\section{\LaTeX 模板的背景与意义}

随着学术规范要求的日益严格，高质量的排版已成为学位论文撰写中不可忽视的一环。相比于传统的文字处理软件，\LaTeX 在处理复杂数学公式、参考文献管理以及长文档结构化排版方面具有显著优势。开发并使用标准化的 \LaTeX 模板，不仅能够显著降低博士生在格式调整上花费的时间成本，使其专注于科研内容的创新，还能确保论文格式的统一性与美观性，提升学术交流的效率与专业度。

\section{研究目标与关键科学问题}

本文的研究目标是面向上海创智学院的博士开题报告撰写需求，设计并实现一套功能完善、易于使用的 \LaTeX 模板。该模板应涵盖开题报告的各个必要组成部分，支持多种学科领域的特殊需求，并提供清晰的使用指南，帮助博士生高效完成开题报告的撰写工作。


\section{全文结构安排}

本文共分为六章，各章内容的安排如下：

第一章：选题背景和意义。介绍了课题的研究背景、研究意义、研究目标及全文结构。

第二章：第一个工作。详细介绍了第一个工作的相关文献、研究方法和初步成果。

第三章：第二个工作。详细介绍了第二个工作的相关文献、研究方法和初步成果。

第四章：第三个工作。详细介绍了第三个工作的相关文献、研究方法和初步成果。

第五章：未来研究计划。总结了已有工作，提出了未来的研究方向和计划。

第六章：学术成果。列出了在课题研究期间发表的论文和取得的专利成果。
	\markboth{基于创智的第一个工作}{基于创智的第一个工作}
\chapter{基于创智的第一个工作}
\label{chap:work1}

\section{第一个工作的相关文献}

第一个工作的相关文献综述内容待补充。

引用一些文章：~\cite{knuth1984tex,latex2025template,wang2024automated,wang2024automated}。

\section{第一个工作的研究方法}

第一个工作的研究方法内容待补充。

\section{第一个工作的初步成果}

第一个工作的初步成果内容待补充。
	\markboth{基于创智的第二个工作}{基于创智的第二个工作}
\chapter{基于创智的第二个工作}
\label{chap:work2}

\section{第二个工作的相关文献}

第二个工作的相关文献内容待补充。

\section{第二个工作的研究方法}

第二个工作的研究方法内容待补充。

\section{第二个工作的初步成果}

第二个工作的初步成果内容待补充。
	\markboth{基于创智的第三个工作}{基于创智的第三个工作}
\chapter{基于创智的第三个工作}
\label{chap:work3}

\section{第三个工作的相关文献}

第三个工作的相关文献内容待补充。

\section{第三个工作的研究方法}

第三个工作的研究方法内容待补充。

\section{第三个工作的初步成果}

第三个工作的初步成果内容待补充。
	\markboth{未来研究计划}{未来研究计划}
\chapter{未来研究计划}
\label{sec:future_intro}

\section{总结与展望}

通过工作一、工作二和工作三，本文设计并实现了一套基于\LaTeX 的上海创智学院博士开题报告撰写模板，解决了研究生在撰写开题报告过程中遇到的格式繁琐、排版困难等问题。目前，模板已经具备了基本的文档结构、样式定制和自动化编译功能，能够满足大多数用户的需求。

然而，本模板仍存在一些不足之处，需要在未来的工作中进一步完善：
\begin{enumerate}
    \item 功能扩展：目前模板主要针对开题报告，未来计划扩展支持学位论文、中期报告等更多类型的文档。
    \item 画质提升：目前的封面Logo清晰度较低，未来计划寻找更高分辨率的图像资源。
\end{enumerate}

\section{研究进度安排}

\begin{table}[H]
  \centering
  \caption{研究进度安排表}
  \begin{tabular}{|p{3cm}|p{7cm}|p{5cm}|}
  \hline
  \multicolumn{1}{|c|}{起讫时间} & \multicolumn{1}{c|}{主要研究内容} & \multicolumn{1}{c|}{预期结果、阶段性成果} \\ \hline
  2023.09 - 2024.05 & 完成工作一。 & 发表论文一。 \\[3ex] \hline
  2024.05 - 2025.05 & 完成工作二。 & 发表论文二。 \\[3ex] \hline
  2025.05 - 2026.05 & 完成工作三。 & 预期发表论文三。 \\[3ex] \hline
  2026.05 - 2027.05 & 完成工作二的扩展。 & 预期发表扩展论文一篇。 \\[3ex] \hline
  2027.05 - 2028.05 & 完成工作三的扩展。 & 预期发表扩展论文一篇。 \\[3ex] \hline
  \end{tabular}
\end{table}

	\chapter{学术创新成果}
\label{chap:plan}
\section{发表论文}

以下作者列表中，“*”表示共同一作。
\begin{enumerate}[label={[\arabic*]}]
    \item \textbf{Xiaozhi Chuang}*, San Zhang*, Si Li. ``A Novel LaTeX Template for Doctoral Research Proposals''. \textit{Journal of Academic Writing}, 2025.
    \item \textbf{Xiaozhi Chuang}, Wu Wang. ``Automated Compilation Systems for Large Scale LaTeX Projects''. \textit{Proceedings of the International Conference on Document Engineering}, 2024.
\end{enumerate}
    
\section{在投论文}
\begin{enumerate}[label={[\arabic*]}]
\item \textbf{Xiaozhi Chuang}, Liu Zhao. ``Optimizing LaTeX Compilation Speed with Parallel Processing''. Submitted to \textit{IEEE Transactions on Software Engineering}.
\end{enumerate}

\section{已授权专利}

\begin{enumerate}[label={[\arabic*]}]
\item 创小智，张三，李四. ``一种基于云端的LaTeX文档协作编辑方法及系统'', ZL 2025 1 0123456.7, 2025年1月1日.
\item 创小智，王五. ``一种LaTeX文档自动化排版与纠错装置'', ZL 2024 1 0987654.3, 2024年12月12日.
\end{enumerate}


\clearpage

    \appendix
    {
    \bibliographystyle{cvpr} % 参考文献使用CVPR会议引用格式
    \bibliography{ref}
    }	
    
	\backmatter

\end{document}

% vim:ts=4:sw=4
